\statement{A.79}{Proved}{A specified finite action forces a noncentral order-four negative image, independently of every frame choice}
There are an actual finite quotient \(Q=\Pi/N_0\),
\(\Pi=C_2\amalg C_3\), and a positive peripheral automorphism
\(h_0\) of Q fixing the literal cusp, with the following properties.
Put
\[
 N_{j+1}=\overline{[N_j,\Pi]},\qquad Q_j=\Pi/N_j.
 \tag{A.79.1}
\]
Every \(Q_j\) is finite. The specified action extends compatibly
to all these actual source-central refinements, with involutory
negative companions and norm-one s- and r-normalizer frames.
One ordinary full positive peripheral automorphism realizes all
these data simultaneously.

Nevertheless, no full positive peripheral h satisfying
\[
                  h(\tau)=\tau,\qquad h_Q=h_0
 \tag{A.79.2}
\]
has \((hc)^2=1\). No s-frame or r-frame is prescribed in this
nonexistence assertion. In fact every such full h preserves a
specified actual finite source quotient \(E_{\rm ord}\) and its
negative companion has exact order four there. The orders are
\[
 |Q|=2^{972}\,72\,7^{13},\qquad
 |E_{\rm ord}|=2^{2700}\,72\,7^{13}.
 \tag{A.79.3}
\]
Thus actual central refinements, even all their compatible
iterations, do not suffice for order-preserving lifting of this
globally negative finite action. The literal-cusp normalization
and the specified finite action remain hypotheses. The later
\(W=0\) actions are not excluded, and the existential
\(\Delta=1\) involution problem remains open.

\begin{proof}
\emph{The specified actual action and the compatible central tests.}\par
Use the changed affine tuple of A.64, before the subsequent
\(W=0\) reselection in A.66. Its exact data are
\path{work_A/A63_zero_t_A56_solution.json}; its SHA256 is
\[
\begin{gathered}
 \texttt{4135fabf1500006d5ec6bbe5dc0571ea}\\
 \texttt{91018abf911e32136cf965729d18d4e9}.
\end{gathered}
 \tag{A.79.4}
\]
The actual finite old model Q has normal forms
\((a,m,q)\), with \(a\in V^8\),
\(m\in\operatorname{Sym}^2_k(k^8)\), \(q\in H\),
\(k=\mathbb F_8\), and \(V=k^{36}\). It uses the exact
opposite-weight group law and C4 source carry of A.56.
The source is the actual pair
\[
 s\longmapsto(d(s),0,H_s),\qquad
 r\longmapsto(0,0,H_r).
 \tag{A.79.5}
\]
In \path{work_A/check_Q6_central_cover.py}, the positive action
is the fixed linear action followed by conjugation by
\texttt{WE}, and the central parity correction
\texttt{PSHIFT}. Substitution of precisely the two entries from
(A.79.4) defines \(h_0\). This specifies an actual finite
automorphism, rather than only a class of conjugators or an
outer action. Canonical c is the actual source automorphism in
the same model. The residual \(P_0=\ker(Q\to H)\) satisfies
\[
              P_0^{\rm ab}=V^8,
 \qquad P_0'=M,
 \qquad M\text{ is H-trivial}.
 \tag{A.79.6}
\]

For the central-refinement assertion, use the actual quotient
\(D=\Pi/U\) of A.65 with the A.64 tuple. Its kernel satisfies
\(U=\overline{[U,\Pi]}\). Its positive action fixes the literal
cusp, its negative companion is involutory, and both quotient
normalizer norms are one. A.65 supplies a full ordinary positive
lift realizing these data. This prior construction is independent
of the ordinary-cup obstruction below.

The finite Q is an action-stable quotient of this D. Since
\(U\subseteq N_0\), induction gives
\[
 U=\overline{[U,\Pi]}\subseteq N_j\quad(j\geq0).
 \tag{A.79.7}
\]
Consequently every actual quotient in (A.79.1) is a quotient of D
and inherits its compatible action and both norm-one frames.
The same full ordinary lift realizes all of them.

Finiteness holds at every step. If \(N_j\) is open, the central
subgroup
\(Z_j=N_j/\overline{[N_j,\Pi]}\) of \(Q_{j+1}\)
has index \(m_j=|Q_j|\) and is topologically finitely generated.
The continuous transfer \(Q_{j+1}\to Z_j\) restricts on this
central subgroup to \(z\mapsto z^{m_j}\). Its image has exponent
dividing six, since \(Q_{j+1}^{\rm ab}\) is a quotient of
\(\Pi^{\rm ab}=C_6\). Thus \(Z_j\) has exponent dividing \(6m_j\),
so it is finite. This proves the induction and, in particular,
the finiteness of the actual first refinement
\(Q_1=\Pi/\overline{[N_0,\Pi]}\).

\emph{The actual ordinary-cup detector.}\par
Let \(a:Q\to V^4\rtimes H\) be the four normalized free
cusp cocycles, writing \(a(g)\) for their additive coordinate.
They satisfy \(a(s)=0\) literally. Put
\(T_0=V^{(3)}\) with its cubed monomial H-action and
\(T=T_0^{16}\). On \(V^4\times T\) use the class-two law
\[
 (a,z)(b,w)=(a+b,z+w+\mathcal C(a,b)),
 \qquad \mathcal C(a,b)_{ij}=a_i^{[2]}*b_j,
 \tag{A.79.8}
\]
for all sixteen ordered pairs \(0\leq i,j\leq3\), with
coordinatewise products and squares. H acts strictly. Pairing
with Q gives the group whose normal forms and product are
\[
 \begin{aligned}
 E_{\rm ord}&=\{(g,z):g\in Q,\ z\in T\},\\
 (g,z)(h,w)&=
 \bigl(gh,\ z+\rho_3(g)w+
       \mathcal C(a(g),\rho(g)a(h))\bigr).
 \end{aligned}
 \tag{A.79.9}
\]
The actual source elements are
\[
 S=(s,0),\qquad
 R=\bigl(r,\mathcal C(a(r),\rho(r)^2a(r))\bigr).
\]
The latter is obtained by conjugating pure r by
\((\rho(r)a(r),0)\) in the ordinary factor, so \(R^3=1\).
This is the source map already used in A.66 and implemented in
\path{work_A/check_A66_ordinary_cubic.py}, with the stored
action (A.79.4).

This is the full actual source image. At a single physical
weight, brackets between distinct free rows give
\((u^2v,v^2u)\), whose binary span is \(k^2\). A trace
annihilator would make a polynomial with distinct powers two
and four vanish on all of k, forcing both coefficients to zero.
The same-row brackets \(u^2v+v^2u\) similarly span k.
The old center pairs opposite physical weights, so these
single-weight brackets have zero old-center component.
They attain all sixteen new copies independently. The already
attained old affine image and the finite Frattini argument of
A.66 therefore give precisely (A.79.9). Its kernel is all T,
which proves the second order formula in (A.79.3).

\emph{Every candidate preserves this detector and has the forced
coefficient action.}\par
The actual normal odd subgroup of H has no invariants on
\(T_0\). Coprime Hochschild--Serre and (A.79.6) give
\[
 \begin{gathered}
 H^1(H,T_0)=H^2(H,T_0)=0,\\
 H^1(Q,T_0)\simeq
       \operatorname{Hom}_H(V^8,T_0)\simeq k^8.
 \end{gathered}
 \tag{A.79.10}
\]
The Hom is over the underlying additive binary modules.
The literal intertwiner is
\(j_3(v)_\lambda=v_{\operatorname{partner}(\lambda)}^4\);
we retain this Frobenius comparison instead of identifying the
two k-representations without it. The eight actual source
classes and 36 coboundaries have rank 44. The source cocycle
space has dimension \(20+24=44\), from the order-two and
order-three source relations, and its coboundary space has
dimension 36. Thus inflation in
(A.79.10) is an isomorphism onto \(H^1(\Pi,T_0)\).
The continuous five-term sequence therefore makes
\[
 \operatorname{Hom}_Q(N_0,T_0)
                 \longrightarrow H^2(Q,T_0)
 \tag{A.79.11}
\]
injective; see \citet{APoonen2017}.

The sixteen extension classes in (A.79.9) are the literal
Frobenius cup products \(e_i^{[2]}\smile e_j\).
They are independent by the actual source attainment above.
The action \(h_0\) fixes the four cusp cohomology classes;
canonical c preserves them with its actual coefficient
intertwiner. Equivariance and (A.79.11) force every full h in
(A.79.2) to preserve all sixteen actual kernel maps and to act
identically on T. Canonical c likewise preserves their joint
kernel and acts by the cubed monomial operator \(C_3\).
Thus descent and the coefficient action are forced for every
candidate; they are not additional restrictions imposed on h.

\emph{Centralizer surjectivity removes every frame choice.}\par
For every \(g\in C_Q(s)\), the literal lift \((g,0)\)
centralizes S in (A.79.9). Indeed both cup terms in the products
with S vanish because \(a(s)=0\), and the old products agree.
Consequently
\[
                  C_{E_{\rm ord}}(S)\longrightarrow C_Q(s)
                   \quad\text{is surjective}.
 \tag{A.79.12}
\]
This includes every centralizing head and every residual affine
or central coordinate.

Let two automorphisms of \(E_{\rm ord}\) cover \(h_0\), act
identically on T, and preserve the conjugacy class of S. Choose
their S-conjugators \(B_1,B_2\). The image of
\(B_1^{-1}B_2\) centralizes s in Q. By (A.79.12), cancel a
lift of this centralizer from \(B_2\); the remaining difference
from \(B_1\) lies in T. Hence the two S-images are T-conjugate.
Their difference cocycle has zero s-restriction class, without
prescribing either conjugator or even its finite head.

\emph{Exhausting all actions over the specified finite action.}\par
Let \(h_*\) be the actual finite action in the saved ordinary
certificate. It fixes \(\tau\) literally and preserves the
S-conjugacy class. Every descended candidate differs from it by
a T-valued 1-cocycle on Q. The preceding paragraph kills its
s-restriction class; the literal cusp kills its tau restriction.
After pullback to the actual source cyclic subgroups, the
complete joint restriction isomorphism is
\[
 H^1(Q,T_0)\longrightarrow
 H^1(\langle s\rangle,T_0)
       \oplus H^1(\langle\tau\rangle,T_0),
 \qquad 8=4+4.
 \tag{A.79.13}
\]
This is the full source restriction of A.56.5, transported by
\(j_3\), also verified by the actual 44/40 coefficient matrices.
The subgroups in (A.79.13) belong to \(\Pi\); no finite cyclic
cohomology group is substituted for the cusp subgroup.

It follows that every candidate has the form
\[
                  h_E=\operatorname{Inn}(v)h_*,
                  \qquad v\in T^\tau.
 \tag{A.79.14}
\]
The vector is unique because \(T^H=0\). The actual negative
coefficient \(C_3\) fixes \(T^\tau\) pointwise. Therefore,
putting \(b_*=h_*c\), we have
\[
 (\operatorname{Inn}(v)b_*)^2
       =\operatorname{Inn}(v+b_*(v))b_*^2=b_*^2.
 \tag{A.79.15}
\]
This exhausts every peripheral action lift of \(h_0\), including
all choices of s- and r-frames.

\emph{The exact nonzero square and verification.}\par
The original-law certificate
\path{work_A/A66_ordinary_cubic_result.json} gives
\[
 b_*^2=\operatorname{Inn}(t),\qquad
 t\in T^\tau,\qquad
                  t_{00,\,\mathrm{weight}\,0}=1\in k.
 \tag{A.79.16}
\]
The equality is checked on the actual S and R. Since \(T^H=0\),
this nonzero t induces a nonidentity inner automorphism.
As T has exponent two, the square in (A.79.16) has order two,
so \(b_*\), and by (A.79.15) every candidate negative image,
has exact order four. A full negative involution would induce
an involution on this preserved finite quotient, a contradiction.

The independent earlier source-law projection is
\path{work_CE/A66_ordinary_projection_independent.json}.
The separate fresh verifier
\path{work_A/code/verify_fixed_action_ordinary_square.py}
preserves the stored witness and checks the actual source and
canonical-c identities, literal cusp, ranks 44/40/36, and all 192
binary cusp-inner directions. It verifies the coordinate one
in (A.79.16) and that the square moves an actual source generator.
Its genuine stdout is the corresponding \path{.out} file;
the complete results and input hashes are in
\path{work_A/fixed_action_ordinary_square_verified.json}.
Its imported legacy group builder performs existing exploratory
algebra but writes no artifacts. No enumeration of the enormous
group in (A.79.3) or search-timeout inference is used.

The finite action (A.79.4) is essential. A.66 and A.75 use
changed actions with zero free cusp inner term and are compatible
with this obstruction. No assertion is made about every finite
action with \(\Delta=1\), an unnormalized full cusp image,
or the universal conjugacy question.
\end{proof}
